
\subsection{A full geometric realization and its complete relator file}

\begin{result}{B.25}{Proved}{All 542 objects and 1620 arrows in a specified realization}
There is an explicitly specified geometric realization of the
nine object families and fourteen named arrow families of
cote 148, scans 31--33, with
\[
 542\text{ objects},\qquad1620\text{ arrows},\qquad
 1079\text{ free cycles},\qquad1080\text{ defining relators}.
                                                               \tag{B.61}
\]
The corresponding path-groupoid functor to $\Pi_1M_{0,5}$
is an equivalence after imposing those relators.  All arrow
images, tree transports and relators are computed in
\texttt{b\_third\_tier.json}, using the exact conventions below
and those of B.22--B.24.  Every arrow has a computed expression
in the sixty core paths and a tree whose edges are twists or
the specified one-node erasures.  The core paths have B.23's
exact four-point covering description.

This theorem specifies the direct $f_{\vec Q,a}$ paths,
the retained-pin erasures, and the three tetrahedral endpoint
choices.  It does not identify a further manuscript homotopy
convention that is not determined by the supplied diagram.
The literal assignment question in B.15 retains precisely
that comparison, together with the general colimit of B.8.
\end{result}
\begin{proof}
Write $\zeta=(-1+i\sqrt3)/2$ and $\xi=(1+i\sqrt3)/2$.
The tetrahedron--triangle model is
\[
 W(q)=(1,\zeta,\zeta^2,-q^{-1},q^{-1}),\qquad
 W=W(\epsilon),\quad\epsilon=1/10.                \tag{B.62}
\]
At $q=0$ its stable limit has a tetrahedral component
$(1,\zeta,\zeta^2,\infty)$ and a three-point tail.
In the coordinate $w=1/z$ near that node, the two
tail marks are $-q,q$.  Thus $wv=q$ is again the normal
plumbing relation.  Exact comparison of the two normalized
cross-ratios for all $120$ permutations shows that $W$
has no nonidentity holomorphic relabelling stabilizer.
This supplies the remaining $120$ objects, indexed by
$g\in\mathfrak S_5$.

The positive sixth-turn is $q(t)=\epsilon e^{\pi it/3}$.
It is homotopic, within the punctured small disc, to
\[
 q(t)=\epsilon(1-t+\xi t).                         \tag{B.63}
\]
Throughout this homotopy the three tetrahedral marks stay
on the unit circle and the two tail marks have modulus
at least $10$, so no collision occurs.  Multiplication
of $W(\epsilon\xi)$ by $\zeta^{-1}$ identifies its endpoint
with $p_\tau W$, where
\[
 p_\tau=(1,2,0,4,3)=(012)(34).                    \tag{B.64}
\]
This is exactly the six-step permutation $\omega\times(-1)$
of the source's seam pins $J\times\varepsilon$.
Indeed the cyclic permutation of $J=\{0,1,2\}$ has
order three and the interchange of $\varepsilon=\{3,4\}$
has order two.  Its six translates concatenate to one
positive turn in the normal coordinate $q$, hence to
the separating meridian, with no suppressed full twist.
The endpoint map for this family is
\[
 \tau_r^0:g\longrightarrow gp_\tau
 \quad(g\in\mathfrak S_5).                        \tag{B.65}
\]

For $\gamma_r^0$ use the path
\[
 E(t)=(1,\zeta,\zeta^2,-(1-t)/\epsilon,
                         1/(\epsilon(1-t))),
 \qquad0\leq t\leq1,                             \tag{B.66}
\]
with the limiting value at $t=1$.  Its two moving marks
stay respectively on $(-\infty,0]$ and $[10,\infty]$,
away from each other and from the tetrahedral triple.
Its endpoint is $(1,\zeta,\zeta^2,0,\infty)$, the
bipyramid already identified with B.20's $K(1/\sqrt3)$.
The source has the conjugate pair $\{1,2\}$ and three
real marks $\{0,3,4\}$ throughout.  After fixing the
three real marks, this component is an upper or lower
half-plane.  Thus (B.66) also records its unique path
class in that real component.  Its endpoint map is
\[
 \gamma_r^0:g\longrightarrow gH_B,
                                                               \tag{B.67}
\]
so it has six incoming arrows at each of the twenty
bipyramids.

It remains to specify the three choices in the
$360$-member family $\gamma_r^{1'}$.  Use the semi-left
model $M$ of (B.54).  Normalize its real labels $0,2,4$
to $0,1,\infty$.  The point with label $1$ becomes
\[
 \lambda_*={-19+i\over181}\in\mathbb H,
\]
and label $3$ becomes its conjugate.  Therefore this
entire real component is the explicit map
\[
 C:\mathbb H\longrightarrow M_{0,5},\qquad
 C(\lambda)=(0,\lambda,1,\overline\lambda,\infty).
                                                               \tag{B.68}
\]
Its inverse takes the labelled cross-ratio determined by
$0,2,4$; thus it is a homeomorphism onto the component,
not merely a dimension count.

Put
\[
 p_0=(0,1,3,4,2),\qquad c=(0\,2\,4),\qquad
 p_j=c^jp_0\quad(j=0,1,2).                         \tag{B.69}
\]
These are the stipulated pin choices.  The tetrahedral
triple at $p_jW$ is the conjugate pair $\{1,3\}$ together
with $c^j(0)$; the other two real marks form its triangular
tail.  Cyclically permuting the three real labels preserves
the half-plane.  The three endpoint coordinates are exactly
\[
\begin{aligned}
 \lambda_0&=-\frac{30}{91}+\frac{110\sqrt3}{819}i,\\
 \lambda_1&=\frac{27}{37}+\frac{30\sqrt3}{407}i,\\
 \lambda_2&=\frac{121}{40}+\frac{33\sqrt3}{40}i.
\end{aligned}                                                     \tag{B.70}
\]
All have positive imaginary part.  Define the representative
arrow of choice $j$ by the straight segment
\[
 C((1-t)\lambda_*+t\lambda_j).                    \tag{B.71}
\]
The upper half-plane is convex, so this is collision free
and gives its unique relative path class there.  Relabelling
by $g$ supplies all arrows, with exact endpoint formula
\[
 \gamma_{r,j}^{1'}:g\ (\mathrm{III}_1)\longrightarrow
                      gp_j\ (\mathrm{II}_1),
 \qquad g\in\mathfrak S_5,\quad j=0,1,2.            \tag{B.72}
\]
There are three distinct targets at each source, and for
each target and each $j$ there is exactly one source.
This proves both valencies three, all $360$ arrows, and
the three free $\mathfrak S_5$-orbits retained in B.12.
The alternative half-plane endpoint pin choices are
not silently equated with (B.69); that selection is the
explicit convention of the present realization.

Choice $j=0$ is supported on a genuine four-point piece.
In (B.68), its three marks $0,\lambda,\overline\lambda$
stay in $|z|<1/2$, because
\[
 |\lambda_*|^2=2/181<1/4,\qquad
 |\lambda_0|^2=400/2457<1/4,
\]
and the disc is convex.  The other two marks are $1$
and $\infty$.  The original radius-$1/4$ circle at the $B$ cluster of $M$
becomes the circle of centre $-8/21$ and radius $20/21$.
Interpolate its centre to $0$ and its radius to $3/4$;
throughout, $0,\lambda_*,\bar\lambda_*$ stay inside
and $1,\infty$ stay outside.  This is an explicit isotopy
of the retained cut.  Keep the circle $|z|=3/4$: its exterior
is a fixed three-point tail, and its interior is the
varying four-point piece with its boundary.  Erasing
the other internal cut gives precisely this one-node
path.  On the retained seam uniformly refine the
two pins to six, using the orientation and the two
existing pins.  This is the uniform overmarking of
cote 148, scan 57, of multiplier three.  It supplies
the target's six pins explicitly; it is not an
identification of two different cardinalities of pin sets.
Thus the $120$ arrows of choice zero are the required
specified operations from a piece of modular dimension one.

Append the $120$ objects in lexicographic permutation
order to B.24.  Append the $120$ sixth-turns, $120$
paths (B.66), then the three $120$-arrow families
(B.71), in the order $j=0,1,2$.  This produces the
$542$ objects and $1620$ arrows of (B.61), retaining
all previous indices.

The path computation uses the same normalization,
basepoint, detour transports and projection as B.21.
For the new paths, choose the following homogeneous representatives:
for (B.63), use $(1:1),(\zeta:1),(\zeta^2:1),(-1:q),(1:q)$;
for (B.66), use $(1:1),(\zeta:1),(\zeta^2:1),
(-(1-t):\epsilon),(1:\epsilon(1-t))$; for (B.71), use
$(0:1),(\lambda:1),(1:1),(\bar\lambda:1),(1:0)$.
Their pair determinants $D_{ij}$ are polynomials of degree
at most two over $\mathbb Q(\sqrt3)(i)$.  On each rational
subinterval, every Bernstein coefficient $z$ of $D_{ij}$
satisfies
\[
 \operatorname{Re}(z\overline{D_{ij}(m)})>0,\qquad
 |\operatorname{Im}(z\overline{D_{ij}(m)})|
 <(2-\sqrt3)\operatorname{Re}(z\overline{D_{ij}(m)}),
\]
where $m$ is the midpoint.  These exact inequalities put
all determinants within $\pi/12$ of their midpoint rays.
The normalized difference of any two finite coordinates is
\[
 \frac{D_{ij}D_{ac}D_{bc}}{D_{ic}D_{jc}D_{ba}},
\]
up to a constant sign, where $a,b,c$ are the three
normalizing labels.  Its argument therefore varies by
less than $6\pi/12=\pi/2$ from its midpoint ray.  This
is the strict half-plane condition for the collision-free
homotopy to its endpoint chord.  A single determinant
certificate consequently proves the homotopy for all
$120$ relabellings.  The five representative families
need respectively $4,8,3,12,16$ subintervals, totaling
$43$.  The code retains their exact Bernstein coefficients
and all $600$ normalized braid images.  Exact crossing
arithmetic and the faithful Artin action verify every
pure-braid word.  An earlier slower run, retained as
\texttt{b\_third\_tier\_bernstein\_wip}, was stopped
before completion; no result uses that incomplete run.

For the spanning tree retain B.24's $421$ edges and
add exactly the $120$ arrows (B.72) with $j=0$.
Their targets run bijectively over the new objects,
so this is a tree with $541$ edges.  Its added edges
are precisely the one-node paths just described.
There are $1620-542+1=1079$ non-tree cycles.  For each
of the $480$ new cycles, the program records its exact
word in B.21's eleven core cycles and imposes that
one relation.  Together with the previous $600$
relations this gives $1080$ relators.  All expanded
words are retained and checked by the Artin action.
Eliminating these $480$ new generators gives exactly
B.24.  This proves completeness of the relator set,
the equivalence to $\Pi_1M_{0,5}$, and the explicit
inverse functor by core words and tree transport.

The same formula is the promised expression for each
arrow through the sixty core paths and the $541$
tree edges: its based cycle is replaced by its recorded
core word, then the source and target tree transports
are restored.  The tree edges are the specified
twists and one-node erasures, and B.23 gives the exact
four-point covering functor for the core.  Thus this
construction includes a computed relative-dimension-one
dictionary for every arrow of the specified full model.
It leaves no unspecified word in that dictionary;
the still-open issue is the explicitly retained comparison
of conventions with the manuscript.
\end{proof}

\begin{result}{B.26}{Proved}{The computed defining-relator counts are minimal}
For the connected graphs used in B.21, B.22, B.24 and
B.25, a groupoid presentation with isotropy
$\Gamma_{0,5}$ requires at least, respectively,
\[
 12,\qquad180,\qquad600,\qquad1080
\]
defining relators.  The presentations constructed there
attain these bounds.
\end{result}
\begin{proof}
The explicit split extension of B.19 is
\[
 1\longrightarrow F(x,y,z)\longrightarrow\Gamma_{0,5}
 \longrightarrow F(a,b)\longrightarrow1.
\]
Each of its six displayed conjugation relations acts
as the identity on the abelianization of the fibre.
The homology spectral sequence has only the columns
and rows numbered $0,1$, since both groups are free.
There is therefore no possible nonzero differential
affecting total degree two, and it gives
\[
 H_1(\Gamma_{0,5};\mathbb Z)=\mathbb Z^5,\qquad
 H_2(\Gamma_{0,5};\mathbb Z)
   =H_1(F_2;\mathbb Z^3)=\mathbb Z^6.              \tag{B.73}
\]
The $H_1$ equality also follows immediately from the
displayed presentation and its split projection.

Collapse a spanning tree of a connected generator graph
and write $b$ for its number of free cycles.  A set of
$r$ defining relators gives a two-dimensional presentation
complex $Y$ with one vertex, $b$ edges and $r$ two-cells.
Its first rational homology has dimension five.  The
cellular chain complex consequently gives
\[
 \dim_{\mathbb Q}H_2(Y;\mathbb Q)=r-b+5.
\]
The classifying map $Y\to B\Gamma_{0,5}$ is surjective
on $H_2$: attach cells of dimensions at least three
to kill the higher homotopy groups of $Y$; these may
quotient $H_2$ but cannot create new degree-two classes.
Equation (B.73) therefore gives $r-b+5\geq6$, or
$r\geq b+1$.  The four cycle counts are $11,179,599,1079$.
Their constructed relator counts are exactly one greater,
which proves the assertion.
\end{proof}
